\documentclass[11pt]{article}
\usepackage[margin=0.78in]{geometry}
\usepackage[T1]{fontenc}
\usepackage{lmodern}
\usepackage{microtype}
\usepackage{xcolor}
\usepackage{enumitem}
\usepackage{longtable}
\usepackage{booktabs}
\usepackage{tabularx}
\usepackage{fancyhdr}
\usepackage{hyperref}
\usepackage{listings}

\definecolor{navy}{RGB}{25,55,92}
\definecolor{pale}{RGB}{239,244,249}
\hypersetup{colorlinks=true,linkcolor=navy,urlcolor=navy}
\setlist{nosep,leftmargin=1.35em}
\setlength{\parindent}{0pt}
\setlength{\parskip}{5pt}
\pagestyle{fancy}
\fancyhf{}
\lhead{Proto-Hive migration plan}
\rhead{Revision 2026-08-12}
\cfoot{\thepage}
\lstset{basicstyle=\ttfamily\small,backgroundcolor=\color{pale},frame=single,
  breaklines=true,columns=fullflexible,showstringspaces=false}

\title{Revised Complete Proto-Hive Migration Plan\\
\large Containerized deployment to ASI replacement VM1 and VM2}
\author{ZeroBot / ProtoCosmo migration workstream}
\date{2026-08-12}

\begin{document}
\maketitle

\begin{center}
\fcolorbox{navy}{pale}{\parbox{0.91\textwidth}{
\textbf{Decision.} Docker Compose is the primary deployment path. The runtime,
supervision, health checks, and receiver configuration live in reproducible
containers. Persistent state and credentials remain in narrowly scoped host
mounts. Cutover is performed one Telegram identity at a time, with exactly one
active receiver and the Pop!\_OS laptop retained as the immediate rollback host.
}}
\end{center}

\section{Purpose and current evidence}

The objective is to migrate four proto-hive identities from the Pop!\_OS laptop
to two existing ASI replacement virtual machines without losing identity,
routing, cursor state, scheduled work, or rollback capability. The target
assignment is:

\begin{center}
\begin{tabular}{lll}
\toprule
Target & Identities & Role \\
\midrule
VM1 & ProtoCosmo, ProtoCosmo2 & Two isolated Compose services \\
VM2 & Protomega, Protomega2 & Two isolated Compose services \\
\bottomrule
\end{tabular}
\end{center}

\textbf{Observed before this revision:}
\begin{itemize}
\item Both targets are reachable Ubuntu 22.04 hosts with 4 CPUs, about 16 GiB
RAM, about 128 GB storage, and clean systemd state.
\item Checksummed, credential-free source bundles have passed integrity,
extraction, compile, and provider-free tests on both targets.
\item Each target passed 113 provider-free tests; one laptop-specific PDF
fixture test was deliberately excluded from remote staging.
\item Fail-closed native staging units and immutable identity roots exist but
remain disabled. No VM receiver has been enabled.
\item A preliminary private ProtoCosmo2 state snapshot is staged on VM1 with
polling disabled. The laptop remains authoritative for every receiver.
\end{itemize}

Containerization changes the packaging and supervision layer, not the safety
contract. Existing verified source pins and state manifests are inputs to the
image and mounted-volume design. The disabled native staging path remains a
fallback until the containerized deployment passes acceptance.

\section{Non-negotiable invariants}

\begin{enumerate}
\item \textbf{Single receiver:} for each bot token, at most one process may poll
Telegram at any instant. A target receiver starts only after its laptop owner
has stopped and absence is independently verified.
\item \textbf{Per-identity isolation:} credentials, cursor/state, sessions,
attachments, logs, and locks are never shared across identities.
\item \textbf{No secrets in artifacts:} image layers, Compose files, Git,
project Markdown, logs, shell arguments, PDF, and chat contain credential
categories or references only, never values.
\item \textbf{Fail closed:} absence of an explicit per-identity activation
marker prevents polling. Health checks cannot enable a receiver.
\item \textbf{Immutable runtime:} source and dependencies are pinned and built
into a versioned image. Runtime containers do not modify their code layer.
\item \textbf{Mutable state is durable:} every required mutable path is a
bounded host bind mount or named volume with verified owner, mode, and backup.
\item \textbf{Rollback before optimization:} the laptop copy is preserved and
restartable through the soak period. No cleanup precedes acceptance.
\item \textbf{One lane at a time:} only one identity enters the stop, final
delta, start, canary, and adjudication sequence at once.
\item \textbf{Evidence before claims:} a lane is complete only when receipts,
receiver counts, fresh canaries, restart checks, and state continuity evidence
exist in the experiment ledger.
\end{enumerate}

\section{Target architecture}

\subsection{Portable container unit}

Build one pinned base image containing the OpenClaw/Node runtime and common
dependencies, plus pinned Omega, PeTTa, and per-identity runner inputs where
needed. Select identity behavior through an explicit service command and
read-only configuration, not through mutation of the image.

Each VM runs a project-scoped Docker Compose stack. Each identity is a separate
service with:
\begin{itemize}
\item its own container name, internal network identity, health check, restart
policy, resource limits, and log policy;
\item read-only runtime/config mounts and separate read-write state mounts;
\item a per-identity secrets directory mounted read-only;
\item no published port unless a verified host integration requires one;
\item polling disabled by default and gated by a host activation marker;
\item a non-root runtime user where repository behavior permits it.
\end{itemize}

Docker Compose provides lifecycle orchestration. Host systemd owns one
Compose-stack unit per VM so the stack starts after Docker and mounted storage,
survives reboot, and has a single inspectable host-level supervisor. Container
restart policy handles process failure; systemd handles Docker/host lifecycle.

\subsection{Host/container boundary}

\begin{longtable}{p{0.25\textwidth}p{0.32\textwidth}p{0.34\textwidth}}
\toprule
Concern & In container/image & On host or mounted storage \\
\midrule
Runtime & Pinned Node/OpenClaw, Python/PeTTa dependencies, reviewed runners &
Docker Engine, Compose plugin, systemd stack unit \\
Identity config & Schema and non-secret defaults & Read-only identity-specific
config and routing policy \\
Credentials & Names and loading mechanism only & Mode-restricted per-identity
secret files or Docker secrets \\
State & State-handling code & Cursor, sessions, memory, schedules, attachment
metadata, durable queues \\
Logs & Structured stdout/stderr contract & Bounded Docker journal/log files and
experiment receipts \\
Health & Local probe executable & Compose health status and host monitoring \\
Receiver activation & Fail-closed check & Per-identity activation marker,
created only during cutover \\
Backups & Restore-compatible formats & Encrypted or access-controlled snapshots
and checksums \\
\bottomrule
\end{longtable}

GitHub CLI authentication and unrelated local developer tooling are explicitly
outside the migration critical path and may be re-established separately by
the operator.

\section{Credential and state categories}

Only categories are listed here. No values, fragments, endpoints containing
secrets, or private-key material belong in this document.

\begin{itemize}
\item VM SSH access and verified host-key records.
\item Per-VM inference API authorization and matching inference user mapping.
\item One Telegram bot token per identity.
\item Internal OpenClaw gateway authorization.
\item Model-provider and Codex authentication profiles.
\item Optional connector OAuth credentials used by a particular identity.
\item Telegram allowlists, chat routing rules, and identity policy files.
\item Persistent state roots, ownership, filesystem modes, and backup location.
\item Optional GitHub/gh authentication, handled outside this migration.
\end{itemize}

The migration copies existing required secrets through bounded, non-logging
mechanisms. Operators manually re-establish only credentials that are
host-bound, expired, unavailable for bounded transfer, or intentionally kept
outside the copied state. Credential rotation remains post-migration cleanup by
current owner decision and is not a cutover gate unless authentication fails or
active misuse is observed.

\section{Implementation phases}

\subsection*{Phase 0: Freeze and reconcile}
\begin{enumerate}
\item Re-audit the four live laptop supervisors and prove exactly one receiver
per identity.
\item Freeze image inputs: repository commits, reviewed dirty-file hashes,
dependency lockfiles, base image digest, and wrapper checksums.
\item Reconcile the current identity delta manifest against every container
mount. Classify each path as immutable image content, read-only config,
read-write state, secret, ephemeral cache, or excluded material.
\item Record disk requirements, identity-to-VM assignment, and rollback owner.
\end{enumerate}

\textbf{Gate 0:} no unclassified production path; no secret in the build
context; all pins and expected hashes recorded.

\subsection*{Phase 1: Build reproducible image and Compose stack}
\begin{enumerate}
\item Create a minimal Dockerfile with a digest-pinned supported base image.
Install dependencies from pinned manifests, copy only allowlisted source, and
run as a non-root user where feasible.
\item Add an entrypoint that validates identity, required mounts, file modes,
state schema, and activation marker before starting. It must exit nonzero on an
unknown identity or incomplete configuration.
\item Define four Compose services across two VM-specific overlays. Use
read-only root filesystems where compatible, temporary filesystems for scratch,
bounded resource/log settings, and no broad host-directory mounts.
\item Add health checks that verify process liveness and local readiness without
calling Telegram or consuming updates.
\item Generate a software bill of materials or at minimum an image digest,
package inventory, source pins, and build transcript.
\end{enumerate}

\textbf{Gate 1:} clean local build; image contains no recognizable secret
material; provider-free test suite passes inside the image; Compose config
renders successfully; all services remain polling-disabled.

\subsection*{Phase 2: Install Docker and stage offline}
\begin{enumerate}
\item Verify official package provenance and install Docker Engine plus Compose
on both existing VMs. No new paid resources are created.
\item Create fixed host directories for images/manifests, per-identity config,
secrets, state, backup snapshots, and receipts. Apply least-privilege modes.
\item Transfer the image by a checksummed archive or pull it from an approved
registry by immutable digest. Transfer Compose manifests separately.
\item Start every service in offline/no-poll mode with placeholder or omitted
Telegram authorization. Run compilation, dependency, state-schema, provider,
and health checks that cannot poll Telegram.
\end{enumerate}

\textbf{Gate 2:} both VMs reproduce hashes; all services become healthy in
offline mode; zero target pollers are observed; laptop pollers remain exactly
one per identity.

\subsection*{Phase 3: Stage private state and rehearse rollback}
\begin{enumerate}
\item Stop each offline service before state import. Copy preliminary state and
required secrets into only that identity's mounts using a bounded allowlist.
\item Verify file counts, hashes where stable, modes, owners, state schema, and
absence of cross-identity files. Do not print content.
\item Start again with polling disabled. Exercise provider authorization and
local response generation without Telegram update consumption where possible.
\item Rehearse container stop, state backup, state restore into a disposable
location, Compose restart, and laptop supervisor restart commands.
\end{enumerate}

\textbf{Gate 3:} private mounts validate; offline provider smoke passes;
backup/restore rehearsal passes; rollback commands and estimated time are
recorded.

\subsection*{Phase 4: Per-identity production cutover}

Use the same transaction for ProtoCosmo2, ProtoCosmo, Protomega, and
Protomega2, one at a time. ProtoCosmo2 remains the first candidate because its
preliminary state is already staged, but the operator may reorder lanes if its
runtime acceptance prerequisites are incomplete.

\begin{enumerate}
\item Announce the lane and create a pre-cutover laptop state snapshot plus
receiver-count receipt.
\item Stop the identity through its owning laptop supervisor. Verify the
process is absent and cannot auto-restart. Leave all other identities running.
\item Copy the final mutable delta, including cursor and pending durable work,
to that identity's target mounts. Verify metadata and state consistency.
\item Create only that identity's activation marker and start its Compose
service. Confirm one target receiver and zero laptop receivers for the token.
\item Send fresh human-authored Telegram canaries: direct/group text as
applicable, reply routing, and a safe attachment. Confirm exactly one response,
correct identity, correct destination, preserved state, and no stale replay.
\item Restart the container and repeat a fresh canary. Inspect bounded logs for
duplicate polling, authentication errors, dropped work, or cross-routing.
\item Accept the lane, or roll it back immediately before touching the next.
\end{enumerate}

\textbf{Gate 4 per identity:} exactly one receiver; all fresh canaries and
restart recovery pass; mutable state advances on the target; rollback remains
available.

\subsection*{Phase 5: Host restart, reboot, and soak}
\begin{enumerate}
\item After all four lanes pass independently, restart each Compose stack under
systemd and verify identity isolation and receiver counts.
\item Reboot one VM at a time. Keep the other VM and laptop rollback artifacts
intact. Confirm Docker, systemd, mounts, and services recover in correct order.
\item Repeat fresh text and attachment canaries after each reboot.
\item Soak for an agreed interval while monitoring process health, restart
counts, disk growth, scheduled work, provider failures, routing, and duplicate
receiver symptoms.
\end{enumerate}

\textbf{Gate 5:} both VMs survive reboot; all four identities pass post-reboot
canaries; soak shows no unresolved severity-one or severity-two defects.

\subsection*{Phase 6: Closeout and deferred cleanup}
\begin{enumerate}
\item Freeze final image digest, Compose files, mount manifest, backup/restore
runbook, test receipts, and identity assignment in the project record.
\item Keep laptop services stopped but restartable until the soak owner accepts
the migration. Archive rather than delete rollback state.
\item Re-establish optional host tooling, including GitHub/gh, outside the
receiver cutover path.
\item Perform deferred credential rotation and revoke superseded access after
migration acceptance. Update secret mounts one identity at a time and retest.
\item Remove the persistent migration cron worker only after all acceptance
gates and the durable record are complete.
\end{enumerate}

\section{Rollback transaction}

Rollback is per identity unless a shared host failure requires reverting both
identities assigned to that VM.

\begin{enumerate}
\item Stop the target Compose service and remove its activation marker.
\item Verify no target receiver remains and prevent automatic restart.
\item Preserve failed-target logs and a checksummed state snapshot for
diagnosis. Do not overwrite the pre-cutover laptop snapshot.
\item If safe and schema-compatible, copy only target state advances needed to
avoid lost durable work; otherwise restore the recorded laptop pre-cutover
state and explicitly record the possible replay window.
\item Start the owning laptop supervisor. Confirm exactly one laptop receiver.
\item Send a fresh canary and verify identity, routing, and state continuity.
\item Record the failure, evidence, and the gate that must pass before retry.
\end{enumerate}

Rollback triggers include duplicate receivers, wrong-chat routing, failed fresh
canaries, cursor regression, repeated crash/restart, mount/permission errors,
unbounded disk growth, or inability to prove receiver exclusivity.

\section{Acceptance matrix}

\begin{longtable}{p{0.23\textwidth}p{0.48\textwidth}p{0.20\textwidth}}
\toprule
Area & Required evidence & Pass condition \\
\midrule
Build & Base digest, image digest, source pins, dependency inventory, clean test
log & Rebuild/test succeeds; no secrets detected \\
Isolation & Compose render, mount inventory, permission audit & No broad or
cross-identity mutable mounts \\
Offline stage & Per-VM hashes, health status, process audit & Healthy services;
zero VM pollers \\
Cutover & Before/after receiver counts and supervisor receipts & Exactly one
receiver for each identity \\
Function & Fresh text, reply-routing, and attachment canaries & Correct single
response in correct chat \\
State & Cursor/session/task continuity receipts & No regression, loss, or
unexpected replay \\
Recovery & Container restart, stack restart, VM reboot receipts & Automatic
recovery with fresh canary pass \\
Rollback & Rehearsal and retained laptop snapshot & Documented command path is
usable within the cutover window \\
Soak & Health, restart, storage, schedule, and routing observations & No open
high-severity defect \\
Closeout & Updated project, experiment, task, decision, and daily memory records
& Evidence is durable and migration worker can retire \\
\bottomrule
\end{longtable}

\section{Operational sequence and ownership}

The migration worker performs one bounded, checkpointed action per cycle. It
must finish within its execution budget, record evidence before the checkpoint
deadline, and never turn an ambiguous state into a cutover. Human operators
provide fresh canary messages and resolve external-account credentials when a
bounded transfer is not appropriate. Ben and Elija have equivalent authority
for IT and permissioning decisions, subject to identity verification and normal
platform safety constraints.

Status reports to Telegram should be short and limited to verified milestones,
actual blockers, a cutover warning, a rollback event, or completion. They must
never include secret values or fragments.

\section{Immediate next actions}

\begin{enumerate}
\item Freeze the Docker/Compose specification and reconcile it against the
existing identity delta manifest.
\item Build the credential-free image locally from the already verified source
pins; run provider-free tests inside it and scan the build context/image for
secret material.
\item Install Docker/Compose on VM1 and VM2 from reviewed official packages,
then transfer the image and manifests by checksum.
\item Pass offline Compose health and state-mount smokes on both targets.
\item Resume the first ProtoCosmo2 lane only after its runtime, mount, and
rollback gates pass; then proceed one identity at a time.
\end{enumerate}

\section{Decision summary}

The containerized approach is adopted because it reduces host drift, makes
runtime pins portable, gives repeatable supervision and health checks, and
simplifies future VM moves. It does not remove the distributed-systems risks
at the Telegram boundary. Single-receiver exclusivity, state transactionality,
fresh external canaries, reboot validation, and laptop rollback therefore
remain the decisive acceptance conditions.

\end{document}
